\section{Analysis of the Numerical Results}
\label{sec:analysis}

This chapter consolidates the per-series numerical results from \cref{sec:classical}, \cref{sec:deep}, and \cref{sec:quantile} into a single comparison and reads off the broader implications. Three concrete questions drove the modelling phase of the project, and we address them in turn.

\subsection{Per-Series Cross-Method Comparison}
\label{sec:analysis:cross_method}

\Cref{tab:cross_method} stacks the best classical, the best deep, and the median-quantile point forecast for each representative series. All three columns use the same weighted skill score on the per-series validation segment.

\begin{table}[h]
\centering
\caption{Best weighted skill score per representative series across the three modelling families. ``Classical'' uses the adaptive 80/20 split from \cref{sec:classical:leaderboard_adaptive}. ``Deep'' uses the architecture results from \cref{tab:deep_results_full}. ``Quantile (median)'' is the point-forecast skill of the LightGBM median quantile from \cref{fig:quantile_panels}.}
\label{tab:cross_method}
\begin{tabular}{l c c c}
\toprule
Series & Classical (best) & Deep (best) & Quantile (median) \\
\midrule
Longest history       & SARIMA $0.0617$        & GRU $0.0686$        & median $0.000$ \\
Highest total weight  & ARIMA(1,1,1) $0.4667$  & LSTM $0.6264$       & \textbf{median $0.865$} \\
Most volatile         & AR(1) $0.8277$         & RNN $0.9235$        & median $0.859$ \\
Most stable           & AR(1) $0.0000$         & all $0.0000$        & median $0.000$ \\
\bottomrule
\end{tabular}
\end{table}

The pattern is clean and tells a single coherent story:

\begin{itemize}
  \item On \textbf{Highest Total Weight}, the most economically important series, the LightGBM median quantile produces the strongest point forecast at $0.865$, comfortably above the deep LSTM at $0.6264$ and the classical ARIMA(1,1,1) at $0.4667$. The series that matters most by weight gets the largest absolute lift from the most flexible model.
  \item On \textbf{Most Volatile}, the deep RNN/GRU and the LightGBM median quantile are essentially tied around $0.86$--$0.92$, both substantially above the classical AR(1) at $0.8277$. Sequence learning and global tree-based learning both find the short-lag persistence; classical AR captures most of it but not all.
  \item On \textbf{Longest History}, every method floors near zero. The series has the most data but the least exploitable autocorrelation, so capacity does not help.
  \item On \textbf{Most Stable}, every method scores $0.0000$. The target itself barely moves, and the weighted skill score does not credit near-zero forecasts on near-zero targets.
\end{itemize}

The headline reading is that on this panel \emph{model choice matters most where the data has structure to exploit}. On the two informative series, moving from classical to deep to quantile-median produces real lifts. On the two near-zero-signal series, no model creates lift, regardless of family.

\subsection{Coverage and the Deployable Artefact}
\label{sec:analysis:coverage}

Point-forecast skill is not the only relevant figure of merit on a hedge-fund panel. An operator scaling positions also needs an interval. \Cref{tab:coverage_summary} summarises the empirical weighted prediction-interval coverage probability (wPICP) and the mean interval width (wMPIW) on each representative series under the rolling chronological backtest of \cref{sec:quantile}.

\begin{table}[h]
\centering
\caption{Empirical coverage and width of the calibrated $80\%$ prediction interval on each representative series. Target wPICP is $0.80$.}
\label{tab:coverage_summary}
\begin{tabular}{l c c c c c}
\toprule
Series & wPICP & PICP & wMPIW & MPIW & Folds \\
\midrule
Longest history       & $0.862$ & $0.860$ & $6.437$       & $6.456$       & $5$ \\
Highest total weight  & $0.792$ & $0.800$ & $0.0005282$   & $0.0005386$   & $2$ \\
Most volatile         & $0.774$ & $0.767$ & $192.1$       & $158.2$       & $5$ \\
Most stable           & $0.762$ & $0.767$ & $0.0003690$   & $0.0003736$   & $3$ \\
\bottomrule
\end{tabular}
\end{table}

Empirical wPICP lands within $\pm 0.04$ of the $0.80$ nominal target on every series, with no systematic over- or under-coverage. The interval widths track the underlying volatility: $192$ on the volatile series, $0.0005$ on the highest-weight series, and $0.0003$ on the most-stable series. That is the calibrated artefact a deployment would consume.

\subsection{Why Smoothing Sometimes Wins}
\label{sec:analysis:smoothing}

A counter-intuitive observation surfaces when comparing the fixed-cutoff classical leaderboard (\cref{tab:classical_leaderboard_fixed}) to the adaptive-split version (\cref{tab:classical_leaderboard_adaptive}). Under the fixed cutoff, smoothing models (SES, Rolling Mean) win three of the four series with skill exactly $0.0000$. Under the adaptive split with longer training windows, more sophisticated models (ARIMA(1,1,1) and AR(1)) take over and the skills become positive.

This is not because smoothing is intrinsically wrong. It is because:

\begin{itemize}
  \item With seven training points (the fixed-cutoff condition on Highest Total Weight and Most Stable), there is not enough data to estimate any likelihood-based model. SES and Rolling Mean are the only family that can be fit at all, so they win by default.
  \item With $124$ to $169$ training points (the adaptive-split condition), likelihood-based ARIMA-family models can be estimated, and on series with real differencing-friendly dynamics they take over. Smoothing's win therefore reverses with data quantity, not with data quality.
\end{itemize}

This makes the methodology choice of moving to adaptive splits substantive rather than cosmetic: it changes which model family wins, in a way that reflects what each family can do when given enough data.

\subsection{Why Recursive Multi-Step Hurts the Highest-Weight Series}
\label{sec:analysis:recursion}

A second observation surfaces from comparing the deep skill on Highest Total Weight ($0.6264$) to the quantile-median skill on the same series ($0.865$). Both models are LightGBM-class learners under similar capacity; the difference is the protocol.

The deep model in \cref{sec:deep} uses recursive multi-step forecasting: the model predicts step $t+1$, then feeds its own prediction back as input for $t+2$, and so on across the validation window. Errors compound across steps, and on a $32$-row validation window with depth-$5$ memory the compounded error grows large enough to cap the skill at $0.6264$.

The quantile model in \cref{sec:quantile} uses one-step quantile prediction with the panel's true features at each row. The model does not feed its own prediction back; it consumes the panel's row-level features each time and produces an interval. This avoids the recursive compounding entirely, and on a series whose row-level features carry signal, it produces the higher $0.865$ skill.

The implication is methodological: \emph{recursive multi-step is the bottleneck on long horizons, not the architecture}. A non-recursive one-step quantile model with feature inputs beats a recursive multi-step LSTM on this series, despite both being roughly comparable in raw learning capacity.

\subsection{Headline Empirical Findings}
\label{sec:analysis:headline}

Bringing the numerical results together:

\begin{itemize}
  \item \textbf{The four representative series cleanly partition into two informative cases (Highest Total Weight, Most Volatile) and two near-zero-signal cases (Longest History, Most Stable).} On the informative cases, modelling produces real lift; on the others, no method beats the floor.
  \item \textbf{The best classical model on the informative cases is ARIMA(1,1,1) on Highest Total Weight ($0.4667$) and AR(1) on Most Volatile ($0.8277$).} Both are simple. Order selection beyond $p,q \le 3$ would not help.
  \item \textbf{The best deep model on the informative cases is LSTM on Highest Total Weight ($0.6264$) and RNN on Most Volatile ($0.9235$).} Sequence learning helps where short-range temporal structure exists.
  \item \textbf{The best one-step quantile-median model on the informative cases is LightGBM on Highest Total Weight ($0.865$) and naive lag-1 on Most Volatile ($0.859$).} The non-recursive protocol is the winner on Highest Total Weight, where the deep recursive forecast compounds error across the validation window.
  \item \textbf{Calibrated $80\%$ prediction intervals are achievable on every representative series}, with empirical wPICP between $0.762$ and $0.862$. That is the deployable artefact a hedge-fund operator would consume.
  \item \textbf{The fixed-cutoff classical study is unfit-table on two of four series}; moving to adaptive 80/20 chronological splits changes which model family wins on three of the four series, validating the methodological change.
\end{itemize}

\subsection{What This Project Does Not Claim}
\label{sec:analysis:limits}

We are explicit about the boundaries of these results:

\begin{itemize}
  \item \textbf{Four series is not the whole panel.} The 36{,}923-series panel is much wider than the four representative cases studied here. The patterns identified in \cref{sec:classical}--\cref{sec:quantile} are likely to generalise to series with similar length, weight, and volatility profiles, but they have not been validated on the full panel.
  \item \textbf{The adaptive split is not directly comparable, slice-for-slice, with the canonical fixed cutoff.} Methodology comparisons between the two are valid; absolute numbers across the two protocols should be read as comparable only within a series.
  \item \textbf{No claim of optimality.} The classical lineup is order-bounded at $p,q \le 3$, the deep lineup is single-layer with hidden size $32$, and the quantile lineup uses one shared hyperparameter configuration. Tighter tuning could move the headline numbers, but probably not their qualitative ranking.
\end{itemize}
